% Conversión editable de PMD-W3-C01-OpenMP.pdf (15 páginas).
% Véase README.md para la correspondencia y las correcciones técnicas.
\documentclass[aspectratio=169,10pt]{beamer}
\usepackage[T1]{fontenc}
\usepackage[utf8]{inputenc}
\usepackage[spanish,es-nodecimaldot,es-noquoting]{babel}
\usepackage{lmodern}
\usepackage{tgheros}
\renewcommand{\familydefault}{\sfdefault}
\usepackage{booktabs,tabularx}
\usepackage{listings}
\usepackage{tikz}
\usetikzlibrary{arrows.meta,positioning,shapes.geometric,calc}

\definecolor{uohblue}{HTML}{0075BA}
\definecolor{ink}{HTML}{162B3A}
\definecolor{muted}{HTML}{536978}
\definecolor{pale}{HTML}{EEF5F9}
\definecolor{codegreen}{HTML}{25784B}
\definecolor{codepurple}{HTML}{7B3979}
\setbeamercolor{normal text}{fg=ink,bg=white}
\setbeamercolor{structure}{fg=uohblue}
\setbeamercolor{frametitle}{fg=ink}
\setbeamercolor{framesubtitle}{fg=uohblue}
\setbeamercolor{block title}{fg=uohblue,bg=pale}
\setbeamercolor{block body}{fg=ink,bg=pale}
\setbeamercolor{alerted text}{fg=uohblue}
\setbeamerfont{frametitle}{size=\LARGE,series=\bfseries}
\setbeamerfont{framesubtitle}{size=\normalsize,series=\mdseries}
\setbeamerfont{block title}{size=\small,series=\bfseries}
\setbeamerfont{block body}{size=\small}
\setbeamertemplate{navigation symbols}{}
\setbeamertemplate{itemize items}[circle]
% Animación progresiva: las listas avanzan elemento a elemento. Las pausas
% explícitas de abajo coordinan código, diagramas y bloques entre columnas.

\setbeamersize{text margin left=8mm,text margin right=8mm}
\setbeamertemplate{frametitle}{%
  \vspace{3mm}\insertframetitle\par
  {\usebeamerfont{framesubtitle}\usebeamercolor[fg]{framesubtitle}\insertframesubtitle\par}
  \vspace{1mm}}
\setbeamertemplate{footline}{%
  \hbox to \paperwidth{\hspace{8mm}\color{muted}\scriptsize
    Alex Di Genova\enspace|\enspace Procesamiento Masivo de Datos
    \hfill OpenMP\qquad\insertframenumber/\inserttotalframenumber\hspace{8mm}}
  \vspace{3mm}}
\setlength{\tabcolsep}{5pt}
\lstdefinestyle{c}{
  language=C,basicstyle=\ttfamily\fontsize{8}{9.7}\selectfont,
  keywordstyle=\color{uohblue}\bfseries,
  commentstyle=\color{codegreen},stringstyle=\color{codepurple},
  backgroundcolor=\color{pale},frame=none,
  columns=fullflexible,keepspaces=true,showstringspaces=false,
  breaklines=true,breakatwhitespace=true,tabsize=4,emptylines=0,
  aboveskip=4pt,belowskip=4pt,xleftmargin=5pt,xrightmargin=5pt,
  framexleftmargin=5pt,framexrightmargin=5pt,
  morekeywords={pragma,omp,parallel,for,sections,section,single,critical,
    atomic,barrier,shared,private,num_threads,schedule,static}}
\lstdefinestyle{shell}{style=c,language=bash,
  basicstyle=\ttfamily\fontsize{8}{10}\selectfont}
\lstset{style=c}
\newcommand{\code}[1]{\texttt{\detokenize{#1}}}
\newcommand{\smallnote}[1]{\par{\scriptsize\color{muted}#1\par}}
\newcommand{\notebookurl}{https://github.com/adigenova/uohpmd/blob/main/code/OpenMP.ipynb}

% Diagramas vectoriales: todos los nodos y las etiquetas son editables.
\tikzset{flow/.style={-{Latex[length=2mm]},thick,draw=uohblue},
  thread/.style={circle,draw=uohblue,fill=pale,minimum size=6mm,
    inner sep=0pt,font=\scriptsize},
  diagramlabel/.style={font=\scriptsize,text=muted,align=center}}
\newcommand{\forkjoin}{%
\begin{tikzpicture}[x=.78cm,y=.64cm]
  \node[diagramlabel] at (0,1.25) {Hilo principal};
  \draw[flow] (0,.95)--(0,.3);
  \draw[thick,uohblue] (-1.5,.3)--(1.5,.3);
  \foreach \x in {-1.5,-.5,.5,1.5}{\draw[flow] (\x,.3)--(\x,-1.1);}
  \draw[thick,uohblue] (-1.5,-1.15)--(1.5,-1.15);
  \node[diagramlabel,right] at (1.7,-.35) {Equipo de hilos};
  \node[diagramlabel,right] at (1.7,-1.15) {Barrera implícita};
  \draw[flow] (0,-1.15)--(0,-2.3);
  \node[diagramlabel,left] at (-.2,-1.75) {Código serial};
  \draw[thick,uohblue] (-1.5,-2.3)--(1.5,-2.3);
  \foreach \x in {-1.5,-.5,.5,1.5}{\draw[flow] (\x,-2.3)--(\x,-3.7);}
  \draw[thick,uohblue] (-1.5,-3.75)--(1.5,-3.75);
  \node[diagramlabel,right] at (1.7,-3) {Equipo de hilos};
  \node[diagramlabel,right] at (1.7,-3.75) {Barrera implícita};
  \draw[flow] (0,-3.75)--(0,-4.55);
\end{tikzpicture}}

\title{Procesamiento Masivo de Datos}
\subtitle{OpenMP}
\author{Alex Di Genova}
\institute{Universidad de O'Higgins}
\date{09/09/2024}
\hypersetup{pdftitle={Procesamiento Masivo de Datos: OpenMP},
  pdfauthor={Alex Di Genova},pdfsubject={Clase introductoria de OpenMP},
  colorlinks=true,urlcolor=uohblue,linkcolor=uohblue}

\begin{document}

% 01 — Original: página 1.
\begin{frame}[plain]
  \begin{tikzpicture}[remember picture,overlay]
    \node[anchor=north west,inner sep=0] at ([xshift=8mm,yshift=-5mm]current page.north west)
      {\includegraphics[height=13mm]{assets/uoh.png}};
    \node[anchor=north east,inner sep=0] at ([xshift=-8mm,yshift=-5mm]current page.north east)
      {\includegraphics[width=49mm]{assets/ici.pdf}};
    \fill[pale] (current page.south west) rectangle ([yshift=14mm]current page.south east);
    \draw[uohblue,line width=1.3pt] ([xshift=8mm,yshift=14mm]current page.south west)
      --([xshift=45mm,yshift=14mm]current page.south west);
  \end{tikzpicture}\par
  \begin{tikzpicture}[remember picture,overlay]
    \node[anchor=north west,inner sep=0,text width=14cm]
      at ([xshift=8mm,yshift=-28mm]current page.north west) {
      {\fontsize{27}{31}\selectfont\bfseries Procesamiento Masivo\\de Datos: \textcolor{uohblue}{OpenMP}\par}
      \vspace{5mm}
      {\large Alex Di Genova\par}
      \vspace{2mm}
      {\color{muted}Universidad de O'Higgins\par}
    };
    \node[anchor=south west,inner sep=0,font=\small]
      at ([xshift=8mm,yshift=7mm]current page.south west) {09/09/2024};
  \end{tikzpicture}
\end{frame}

% 02 — Original: página 2.
\begin{frame}
  \vfill
  {\fontsize{36}{40}\selectfont\bfseries OpenMP\par}
  \par\vspace{4mm}
  {\Large\color{uohblue}Introducción y funciones\par}
  \par\vspace{7mm}
  {\color{muted}Memoria compartida · Regiones paralelas · Trabajo compartido\par}
  \vfill
\end{frame}

% 03 — Original: página 3.
\begin{frame}[fragile]{OpenMP}{Introducción}
\small
\begin{columns}[T,onlytextwidth]
\column{.54\textwidth}
  \begin{itemize}[<+->]\setlength\itemsep{4pt}
    \item API para programar sistemas de \alert{memoria compartida} en C, C++ y Fortran.
    \item Combina directivas del compilador, funciones de tiempo de ejecución y variables de entorno.
    \item El programador identifica el paralelismo e inserta las directivas apropiadas.
    \item En C/C++, las directivas comienzan con \code{#pragma omp}.
  \end{itemize}
\column{.42\textwidth}
  \begin{block}{Sintaxis de una región paralela}
\begin{lstlisting}
#pragma omp parallel [clausulas]
{
    /* Bloque estructurado */
}
\end{lstlisting}
  \end{block}
  \smallnote{Los corchetes indican elementos opcionales; no son parte del código C.}
  \par\vspace{3mm}
  El programador decide:
  \begin{itemize}[<+->]\setlength\itemsep{4pt}
    \item qué código se ejecuta en paralelo;
    \item qué variables son privadas o compartidas;
    \item cómo sincronizar los accesos.
  \end{itemize}
\end{columns}
\end{frame}

% 04 — Original: página 4.
\begin{frame}{OpenMP}{La región paralela · Modelo fork--join}
\small
\begin{columns}[c,onlytextwidth]
\column{.58\textwidth}
  \begin{itemize}[<+->]\setlength\itemsep{2pt}
    \item La ejecución comienza en un hilo inicial. Una región \code{parallel} forma un \alert{equipo de hilos}.
    \item El hilo que encuentra la región es el hilo principal del equipo.
    \item Todos ejecutan el bloque. Las construcciones de trabajo compartido distribuyen tareas.
    \item Al terminar la región hay una \alert{barrera implícita}; continúa el hilo que la encontró.
    \item Evitar regiones pequeñas ayuda a reducir el sobrecosto.
  \end{itemize}
  \begin{block}{Número de hilos solicitado}
    \small\code{OMP_NUM_THREADS} · \code{omp_set_num_threads()}\\[2pt]
    Cláusula \code{num_threads(n)} en una región concreta.
  \end{block}
\column{.39\textwidth}
  \centering\forkjoin
\end{columns}
\end{frame}

% 05 — Original: página 5. Figura reconstruida en TikZ.
\begin{frame}{OpenMP}{Modelo de memoria}
\begin{columns}[c,onlytextwidth]
\column{.49\textwidth}
  \begin{block}{Variables privadas}
    Cada hilo dispone de su propia copia. Un mismo nombre puede referirse a ubicaciones distintas de memoria.
  \end{block}
  \begin{block}{Variables compartidas}
    Los hilos pueden leer y escribir la misma ubicación. El programador debe coordinar los accesos que puedan entrar en conflicto.
  \end{block}
  \small Las variables globales se comparten por defecto, salvo excepciones como \code{threadprivate}.
  \par\vspace{2mm}
  \smallnote{Privatizar un puntero no privatiza automáticamente los datos a los que apunta.}
\column{.48\textwidth}
  \centering
  \begin{tikzpicture}[x=1cm,y=1cm]
    \node[ellipse,fill=pale,draw=uohblue,thick,align=center,
      minimum width=2.5cm,minimum height=1.25cm] (mem) at (0,0) {Memoria\\compartida};
    \foreach \ang/\id in {90/0,30/1,-30/2,-90/3,-150/4,150/5}{
      \node[thread] (t\id) at (\ang:1.7) {H\id};
      \node[ellipse,draw=muted,fill=white,font=\scriptsize,align=center,
        inner sep=3pt] (p\id) at (\ang:2.65) {Memoria\\privada};
      \draw[thick,uohblue] (mem)--(t\id);
      \draw[muted] (t\id)--(p\id);
    }
  \end{tikzpicture}\par
\end{columns}
\end{frame}

% 06 — Original: página 6.
\begin{frame}[fragile]{OpenMP}{Construcciones de trabajo compartido}
\small
  Dentro de una región paralela, estas construcciones distribuyen trabajo entre los hilos del equipo.
  \par\vspace{3mm}
\begin{columns}[T,onlytextwidth]
\column{.31\textwidth}
  \begin{block}{\code{for}: iteraciones}
    \small Reparte las iteraciones de un ciclo entre los hilos.
  \end{block}
\begin{lstlisting}
#pragma omp for
for (int i = 0; i < n; i++) {
    procesar(i);
}
\end{lstlisting}
\column{.34\textwidth}
  \begin{block}{\code{sections}: bloques}
    \small Cada sección se ejecuta una vez, por uno de los hilos.
  \end{block}
\begin{lstlisting}
#pragma omp sections
{
    #pragma omp section
    { funcion_a(); }

    #pragma omp section
    { funcion_b(); }
}
\end{lstlisting}
\column{.30\textwidth}
  \begin{block}{\code{single}: un hilo}
    \small Un solo hilo ejecuta el bloque; no se especifica cuál.
  \end{block}
\begin{lstlisting}
#pragma omp single
{
    preparar_datos();
}
\end{lstlisting}
\end{columns}
  \par\vspace{3mm}
  \smallnote{Por defecto hay una barrera al final de estas construcciones. La cláusula \code{nowait} permite omitirla cuando sea correcto. Fragmentos ilustrativos.}
\end{frame}

% 07 — Original: página 7.
\begin{frame}[fragile]{OpenMP}{Paralelismo anidado}
\small
\begin{columns}[c,onlytextwidth]
\column{.48\textwidth}
  Una región paralela puede aparecer dentro de otra. Cada hilo que la encuentra forma un nuevo equipo.
  \par\vspace{3mm}
  \begin{tikzpicture}[x=1cm,y=.8cm]
    \node[diagramlabel] at (0,.7) {Hilo inicial};
    \draw[flow] (0,.4)--(0,0);
    \draw[uohblue,thick] (-1.3,0)--(1.3,0);
    \foreach \x in {-1.3,1.3}{
      \draw[flow] (\x,0)--(\x,-.9);
      \draw[uohblue,thick] (\x-.55,-.9)--(\x+.55,-.9);
      \draw[flow] (\x-.55,-.9)--(\x-.55,-2.15);
      \draw[flow] (\x+.55,-.9)--(\x+.55,-2.15);
      \draw[uohblue,thick] (\x-.55,-2.2)--(\x+.55,-2.2);
      \draw[flow] (\x,-2.2)--(\x,-3.1);
    }
    \draw[uohblue,thick] (-1.3,-3.15)--(1.3,-3.15);
    \draw[flow] (0,-3.15)--(0,-3.7);
    \node[diagramlabel,right] at (2,-1.5) {Equipos\\internos};
  \end{tikzpicture}\par
  \smallnote{Esquema: 2 hilos externos y 2 hilos por equipo interno, si la configuración permite ambos niveles activos.}
\column{.48\textwidth}
  \textbf{Ejemplo del original: 1 $\times$ 5}
\begin{lstlisting}
#pragma omp parallel num_threads(1)
{
    Work1();

    #pragma omp parallel num_threads(5)
    {
        Work2();
    }
}
\end{lstlisting}
  \small La región externa tiene un hilo; la interna solicita cinco.
  \par\vspace{3mm}
  \begin{block}{Niveles activos}
    \small Para dos niveles con varios hilos: \code{omp_set_max_active_levels(2)}. El entorno puede limitar el tamaño de los equipos.
  \end{block}
\end{columns}
\end{frame}

% 08 — Original: página 8.
\begin{frame}[fragile]{OpenMP}{Sincronización}
\small
\begin{columns}[T,onlytextwidth]
\column{.31\textwidth}
  \begin{block}{Barrera}
    Todos los hilos del equipo esperan hasta que todos hayan llegado al mismo punto.
  \end{block}
\begin{lstlisting}
producir();
#pragma omp barrier
consumir();
\end{lstlisting}
  \small Debe ser alcanzada por todos los hilos del equipo.
\column{.31\textwidth}
  \begin{block}{Región crítica}
    Solo un hilo a la vez ejecuta una región crítica con el mismo nombre.
  \end{block}
\begin{lstlisting}
#pragma omp critical
{
    sum += c;
    contador++;
}
\end{lstlisting}
  \small Protege un bloque de instrucciones. Las regiones sin nombre comparten la misma exclusión.
\column{.31\textwidth}
  \begin{block}{Operación atómica}
    Protege una operación admitida sobre una ubicación concreta de memoria.
  \end{block}
\begin{lstlisting}
#pragma omp atomic update
sum += c;
\end{lstlisting}
  \small No protege un bloque arbitrario de código.
\end{columns}
\par\vspace{4mm}
\smallnote{\code{critical} y \code{atomic} no son barreras. Los accesos concurrentes a un dato compartido deben usar una sincronización compatible.}
\end{frame}

% 09 — Original: página 9.
\begin{frame}{OpenMP}{Funciones de tiempo de ejecución}
  \small Permiten consultar y modificar la configuración. Por ejemplo, \code{OMP_NUM_THREADS} fija una solicitud inicial y \code{omp_set_num_threads()} la cambia para regiones posteriores del contexto de la tarea que llama.
  \par\vspace{2mm}
  \footnotesize
  \renewcommand{\arraystretch}{1.12}
  \begin{tabularx}{\textwidth}{@{}lX@{}}
  \toprule
  \textbf{Función} & \textbf{Descripción} \\
  \midrule
  \code{omp_set_num_threads(n)} & Solicita el número de hilos para regiones posteriores. \\
  \code{omp_get_num_threads()} & Número de hilos del equipo actual; fuera de él, 1. \\
  \code{omp_get_max_threads()} & Cota superior para un nuevo equipo sin cláusula \code{num_threads}. \\
  \code{omp_get_num_procs()} & Número de procesadores disponibles para el dispositivo. \\
  \code{omp_get_thread_num()} & Identificador del hilo: desde 0 hasta el tamaño del equipo menos 1. \\
  \code{omp_in_parallel()} & Indica si se ejecuta dentro de una región paralela activa. \\
  \code{omp_get_dynamic()} & Consulta si el ajuste dinámico del número de hilos está habilitado. \\
  \code{omp_set_dynamic(valor)} & Habilita o deshabilita ese ajuste. \\
  \code{omp_get_nested()} & Consulta el paralelismo anidado; función obsoleta. \\
  \code{omp_set_nested(valor)} & Configura el paralelismo anidado; función obsoleta. \\
  \bottomrule
  \end{tabularx}
  \par\vspace{2mm}
  \smallnote{Se conservan las dos últimas funciones del original. Alternativas: \code{omp_get_max_active_levels()} y \code{omp_set_max_active_levels(n)}.}
\end{frame}

% 10 — Original: página 10.
\begin{frame}[fragile]{Ejemplos OpenMP}{Hello world}
\begin{columns}[T,onlytextwidth]
\column{.55\textwidth}
  \lstinputlisting{examples/hello.c}
  \par\vspace{2mm}
  \textbf{Compilar y ejecutar en una terminal}
\begin{lstlisting}[style=shell]
gcc -fopenmp examples/hello.c -o hello
export OMP_NUM_THREADS=3
./hello
\end{lstlisting}
\column{.41\textwidth}
  \begin{block}{Un saludo por hilo}
    Cada hilo ejecuta el mismo bloque y consulta su identificador con \code{omp_get_thread_num()}.
  \end{block}
  \par\vspace{2mm}
  \textbf{Una salida posible}
\begin{lstlisting}[language={}]
Hola Mundo... desde hilo = 1
Hola Mundo... desde hilo = 0
Hola Mundo... desde hilo = 2
\end{lstlisting}
  \small El orden de los mensajes puede cambiar entre ejecuciones.
  \par\vspace{3mm}
  \smallnote{En Jupyter/Colab: \texttt{\%env OMP\_NUM\_THREADS=3}; anteponer \code{!} a los comandos de compilación y ejecución.}
\end{columns}
\end{frame}

% 11 — Original: página 11 (primer ejemplo).
\begin{frame}[fragile]{Ejemplos OpenMP}{Un for dentro de parallel: trabajo repetido}
\begin{columns}[T,onlytextwidth]
\column{.54\textwidth}
  \lstinputlisting{examples/for_replicado.c}
\begin{lstlisting}[style=shell]
gcc -fopenmp examples/for_replicado.c -o for1
OMP_NUM_THREADS=3 ./for1
\end{lstlisting}
\column{.42\textwidth}
  \begin{block}{Cada hilo recorre todo el ciclo}
    Con un equipo de 3 hilos, cada hilo ejecuta 10 iteraciones: \alert{30 mensajes en total}.
  \end{block}
  \par\vspace{3mm}
  \begin{tikzpicture}[x=.39cm,y=.7cm]
    \foreach \h in {0,1,2}{
      \node[diagramlabel,left] at (-.4,-\h) {H\h};
      \foreach \k in {0,...,9}{
        \node[draw=uohblue,fill=pale,minimum width=3.7mm,
          minimum height=4.5mm,inner sep=0pt,font=\tiny] at (\k,-\h) {\k};
      }
    }
  \end{tikzpicture}\par
  \par\vspace{3mm}
  \smallnote{Declarar \code{k} dentro del bloque evita que los hilos compartan el índice y compitan al modificarlo. \code{parallel} por sí solo no reparte el ciclo.}
\end{columns}
\end{frame}

% 12 — Original: página 11 (segundo ejemplo).
\begin{frame}[fragile]{Ejemplos OpenMP}{omp for: distribuir las iteraciones}
\begin{columns}[T,onlytextwidth]
\column{.54\textwidth}
  \lstinputlisting{examples/for_repartido.c}
\begin{lstlisting}[style=shell]
gcc -fopenmp examples/for_repartido.c -o for2
OMP_NUM_THREADS=2 ./for2
\end{lstlisting}
\column{.42\textwidth}
  \begin{block}{Cada iteración se ejecuta una vez}
    Con 2 hilos y \code{schedule(static)}, las 10 iteraciones se reparten en dos bloques de 5.
  \end{block}
  \par\vspace{3mm}
  \begin{tikzpicture}[x=.43cm,y=.8cm]
    \node[diagramlabel,left] at (-.4,0) {H0};
    \node[diagramlabel,left] at (-.4,-1) {H1};
    \foreach \k in {0,...,4}{
      \node[draw=uohblue,fill=pale,minimum width=4mm,minimum height=5mm,
        inner sep=0pt,font=\tiny] at (\k,0) {\k};
    }
    \foreach \k in {5,...,9}{
      \node[draw=uohblue,fill=pale,minimum width=4mm,minimum height=5mm,
        inner sep=0pt,font=\tiny] at (\k,-1) {\k};
    }
  \end{tikzpicture}\par
  \par\vspace{3mm}
  \small El índice del ciclo de \code{omp for} es privado. Hay una barrera implícita al final del ciclo.
  \par\vspace{2mm}
  \smallnote{También se puede escribir \code{#pragma omp parallel for}. El orden de impresión sigue siendo no determinista.}
\end{columns}
\end{frame}

% 13 — Original: página 12 (funciones auxiliares).
\begin{frame}[fragile]{Ejemplos OpenMP}{Funciones de trabajo para sections, single y anidamiento}
  \small Las funciones de \code{work.h} imprimen el trabajo y el identificador del hilo.
  \par\vspace{2mm}
\begin{columns}[T,onlytextwidth]
\column{.48\textwidth}
  \lstinputlisting[firstline=3,lastline=5]{examples/work.h}
  \lstinputlisting[firstline=7,lastline=14,basicstyle=\ttfamily\fontsize{7}{8.5}\selectfont]{examples/work.h}
\column{.48\textwidth}
  \lstinputlisting[firstline=15,lastline=21,basicstyle=\ttfamily\fontsize{7}{8.5}\selectfont]{examples/work.h}
  \begin{block}{Simular trabajo}
    \small \code{sleep(1)} espera un segundo en \code{Work1}, \code{Work2} y \code{Work4}.
  \end{block}
\end{columns}
  \par\vspace{3mm}
  \smallnote{Archivo auxiliar compartido por los siguientes ejemplos. \code{unistd.h} y \code{sleep} requieren un entorno POSIX, como Linux o Colab.}
\end{frame}

% 14 — Original: página 12 (sections).
\begin{frame}[fragile]{Ejemplos OpenMP}{Secciones paralelas}
\begin{columns}[T,onlytextwidth]
\column{.48\textwidth}
  \lstinputlisting{examples/sections.c}
\column{.48\textwidth}
  \begin{itemize}[<+->]\setlength\itemsep{5pt}
    \item Tres secciones para cinco hilos: algunos no reciben trabajo.
    \item \code{Work2()} y \code{Work3()} pertenecen a la misma sección: se ejecutan en orden y en el mismo hilo.
  \end{itemize}
\begin{lstlisting}[style=shell]
gcc -fopenmp examples/sections.c -o sections
OMP_NUM_THREADS=5 ./sections
\end{lstlisting}
  \textbf{Una salida posible}
\begin{lstlisting}[language={}]
executing work 1 hilo:1
executing work 2 hilo:2
executing work 4 hilo:3
executing work 3 hilo:2
\end{lstlisting}
  \smallnote{El reparto de secciones entre hilos puede variar.}
\end{columns}
\end{frame}

% 15 — Original: página 12 (single).
\begin{frame}[fragile]{Ejemplos OpenMP}{single: un bloque ejecutado por un hilo}
\begin{columns}[T,onlytextwidth]
\column{.48\textwidth}
  \lstinputlisting{examples/single.c}
\begin{lstlisting}[style=shell]
gcc -fopenmp examples/single.c -o single
OMP_NUM_THREADS=10 ./single
\end{lstlisting}
\column{.48\textwidth}
  \begin{block}{Con un equipo de 10 hilos}
  \begin{tabularx}{\linewidth}{@{}l>{\raggedright\arraybackslash}X@{}}
    \code{Work1()} & 10 ejecuciones: una por hilo.\\[5pt]
    \code{Work2()} & 1 ejecución.\\[5pt]
    \code{Work3()} & 1 ejecución, en el mismo hilo que \code{Work2()}.\\[5pt]
    \code{Work4()} & 10 ejecuciones, después de la barrera de \code{single}.
  \end{tabularx}
  \end{block}
  \par\vspace{3mm}
  El hilo que ejecuta \code{single} no tiene por qué ser el hilo principal.
  \par\vspace{3mm}
  \smallnote{No hay barrera de entrada: \code{Work2()} puede comenzar mientras otros hilos aún ejecutan \code{Work1()}. Todos esperan al final de \code{single}.}
\end{columns}
\end{frame}

% 16 — Original: página 13 (anidamiento).
\begin{frame}[fragile]{Ejemplos OpenMP}{Regiones paralelas anidadas}
\begin{columns}[T,onlytextwidth]
\column{.48\textwidth}
  \lstinputlisting{examples/nested.c}
\begin{lstlisting}[style=shell]
gcc -fopenmp examples/nested.c -o nested
OMP_DYNAMIC=FALSE ./nested
\end{lstlisting}
\column{.48\textwidth}
  \textbf{Una salida posible}
\begin{lstlisting}[language={}]
executing work 1 hilo:0
executing work 2 hilo:2
executing work 2 hilo:4
executing work 2 hilo:0
executing work 2 hilo:1
executing work 2 hilo:3
\end{lstlisting}
  \begin{itemize}[<+->]\setlength\itemsep{5pt}
    \item El equipo externo tiene un único hilo.
    \item La región interna solicita un equipo de cinco hilos.
    \item Los identificadores son locales a cada equipo.
  \end{itemize}
  \par\vspace{2mm}
  \smallnote{La región externa de un hilo es inactiva: este caso solo necesita un nivel activo.}
\end{columns}
\end{frame}

% 17 — Original: página 13 (critical).
\begin{frame}[fragile]{Ejemplos OpenMP}{Acumulación compartida con una sección crítica}
\begin{columns}[T,onlytextwidth]
\column{.60\textwidth}
  \lstinputlisting[basicstyle=\ttfamily\fontsize{7.6}{9}\selectfont]{examples/critical.c}
\column{.36\textwidth}
  \begin{block}{Proteger la actualización}
    \small \code{sum += c} lee y escribe la variable compartida. \code{critical} permite una actualización a la vez.
  \end{block}
\begin{lstlisting}[style=shell]
gcc -fopenmp \
    examples/critical.c -o critical
OMP_NUM_THREADS=5 ./critical
\end{lstlisting}
  \small Cada iteración imprime su contribución; al salir del ciclo se muestra la suma total.
  \par\vspace{3mm}
  \smallnote{\code{rand()} se llama en serie. Sus valores dependen de la implementación de C.}
\end{columns}
\end{frame}

% 18 — Original: página 14. Captura original extraída del PDF.
\begin{frame}{Google Colab}{Material práctico de la clase}
\begin{columns}[c,onlytextwidth]
\column{.40\textwidth}
  \centering\includegraphics[height=5.6cm]{assets/colab.png}
\column{.56\textwidth}
  {\Large\bfseries Notebook de OpenMP\par}
  \par\vspace{4mm}
  Ejemplos de compilación y ejecución en un notebook.
  \par\vspace{5mm}
  \href{\notebookurl}{\beamergotobutton{Abrir OpenMP.ipynb en GitHub}}
  \par\vspace{4mm}
  {\footnotesize\url{https://github.com/adigenova/uohpmd/blob/main/code/OpenMP.ipynb}\par}
  \par\vspace{4mm}
  \smallnote{Enlace y captura conservados del PDF original.}
\end{columns}
\end{frame}

% 19 — Original: página 15.
\begin{frame}{¿Consultas?}
  \vfill
  \begin{center}
    {\LARGE ¿Consultas o comentarios?\par}
    \par\vspace{8mm}
    {\Large\color{uohblue}Muchas gracias\par}
  \end{center}
  \vfill
\end{frame}

\end{document}
